% Partie : méthodes
% Chapitre : démonstrations
% Section : je n'y arrive pas

\documentclass[../../../include/open-logic-section]{subfiles}

\begin{document}
\olfileid[fr]{mth}{prf}{cnt}
\olsection{Je n'y arrive pas !{}}

Nous connaissons tous des moments où
nous avons envie d'abandonner. Pourtant,
vous \emph{pouvez} y arriver.
Votre enseignant ou enseignante,
la personne qui l'assiste et vos
camarades peuvent vous aider :
demandez-leur conseil !{} Voici
quelques moyens d'éviter d'en arriver
à une crise et quelques idées pour
les moments où vous vous sentez
prêt à renoncer.

Pour mettre toutes les chances
de votre côté lorsque vous
résolvez des problèmes :
\begin{enumerate}
\item \emph{Commencez le plus tôt possible.}
  Le semestre est chargé et beaucoup
  d'entre nous ont tendance à remettre
  le travail au lendemain. L'une des
  meilleures choses à faire est donc
  de commencer les devoirs tôt. Si
  vous êtes bloqué, vous aurez ainsi
  le temps de chercher une solution
  autre que de pleurer.
\item \emph{Parlez avec vos camarades.}
  Vous n'êtes pas seul : d'autres
  personnes du cours rencontrent
  peut-être des difficultés, pas
  forcément les mêmes que les vôtres.
  En discuter avec elles peut vous
  offrir une autre perspective et
  débloquer votre réflexion. Ne
  recopiez pas leur solution :
  demandez plutôt un indice,
  expliquez où vous butez et
  demandez quelle pourrait être
  l'étape suivante. Lorsque vous
  aurez compris, rendez la pareille.
  Aider quelqu'un et lui expliquer
  une idée vous aidera également
  à mieux la comprendre.
\item \emph{Demandez de l'aide.}
  De nombreuses ressources sont
  à votre disposition. Votre
  enseignant ou enseignante et
  la personne qui l'assiste
  \emph{souhaitent} votre réussite.
  Ces personnes peuvent vous aider à
  analyser un problème et à
  repérer l'étape qui vous
  met en difficulté.
\item \emph{Faites une pause.}
  Si vous êtes bloqué, c'est
  \emph{peut-être} parce que
  vous regardez ce problème
  depuis trop longtemps.
  Accordez-vous une courte
  pause, prenez un thé ou
  travaillez un moment sur
  un autre problème, puis
  revenez-y l'esprit reposé.
  Une nuit de sommeil peut
  aussi aider.
\end{enumerate}

Remarquez que ces méthodes supposent
que vous avez commencé la démonstration
bien à l'avance. Si vous la commencez
à deux heures du matin, la veille
de la date limite, elles risquent
de moins vous aider.

Cela peut paraître décourageant,
mais résoudre un problème de
démonstration est un défi
gratifiant. Certaines personnes
en font même leur métier :
elles doivent donc y trouver
quelque plaisir. Comme presque
toute activité, la résolution
de problèmes et la pratique
des démonstrations demandent
de l'entraînement. Si des
camarades semblent y arriver
facilement, c'est probablement
qu'ils se sont déjà beaucoup
exercés. Ne renoncez pas trop vite.

Si le temps ou la patience vous
manque pour un problème donné,
ce n'est pas grave. Cela ne
signifie ni que vous manquez
d'intelligence ni que vous
n'y arriverez jamais. Demandez
à votre enseignant, à votre
enseignante ou à une autre
personne du cours comment le résoudre,
puis repérez l'endroit où
vous vous êtes trompé ou
êtes resté bloqué. Vous
éviterez ainsi le même
obstacle la prochaine fois.
Essayez ensuite de refaire
le problème sans regarder
la solution. Et la prochaine
fois, commencez plus tôt
et demandez de l'aide
plus tôt.

\end{document}
